\section{Physical charge lattice and denominator polarization}
\label{sec:physical-charge-lattice-duality}

The compactification supplies a charge problem before one chooses a
Siegel cusp.  The cusp product is obtained after three reductions:
pass from microscopic charges to their rank-two Gram invariants, choose
the standard Fock polarization, and retain only the active Jacobi
support.  The Lorentzian denominator lattice then appears by an explicit
anti-isometry.  The bracket-level BPS algebra is not part of this
derivation.

\begin{assumption}[Physical two-charge input]
The protected sector is indexed, at the level of the index, by pairs of
charges $(Q,P)$ in an even integral charge lattice.  The protected
index depends on the pair only through the three T-duality invariants
$$
n=\frac{(Q,Q)}2,\qquad l=(Q,P),\qquad m=\frac{(P,P)}2.
$$
Equivalently, it factors through the half-integral symmetric matrix
$$
T(\gamma)=T(n,l,m)=
\begin{pmatrix}n&l/2\\ l/2&m\end{pmatrix},
\qquad \gamma=(n,l,m)\in\Z^3.
$$
This is the Harvey--Moore/Dijkgraaf--Verlinde--Verlinde charge
grading input \cite{HM,DVV}.  It is physical input, not a theorem
proved from the curve-counting geometry in this paper.
\end{assumption}

\begin{proposition}[Gram-invariant index lattice]
After passing to rank-two Gram invariants, the universal index lattice is
$$
\Gamma_{\mathrm{BPS}}=\Z^3
=\{(n,l,m)\mid T(n,l,m)\in\operatorname{Sym}_2(\tfrac12\Z),
\ T_{11},T_{22}\in\Z,\ 2T_{12}\in\Z\}.
$$
Its discriminant is
$$
\Delta(\gamma)=\det(2T(\gamma))=4nm-l^2.
$$
The polarized bilinear form is
$$
\langle\gamma,\gamma'\rangle_{\mathrm{BPS}}
=2(nm'+n'm)-ll',
\qquad \gamma'=(n',l',m').
$$
\end{proposition}

\begin{proof}
Evenness of the charge lattice gives
$(Q,Q)/2,(P,P)/2\in\Z$ and $(Q,P)\in\Z$.  Hence the Gram matrix of
$(Q,P)$ maps to $2T(n,l,m)$ in the displayed lattice.  The Borcherds
product uses the full automorphic completion of this Gram-invariant
lattice; no assertion is made here that every triple is realized by a
microscopic pair in a fixed compactification.  The determinant is
$4nm-l^2$.  The displayed pairing is one half of the polarization of
this quadratic form:
$$
\Delta(\gamma+\gamma')-\Delta(\gamma)-\Delta(\gamma')
=2\langle\gamma,\gamma'\rangle_{\mathrm{BPS}}.
$$
\end{proof}

\begin{proposition}[Anti-isometry to the type-II Lorentzian lattice]
Let
$$
\La^{2,1}_{II}=2\Z f_2\oplus\Z f_3\oplus 2\Z f_{-2},
\qquad
(f_2,f_{-2})=-1,\qquad (f_3,f_3)=2.
$$
Define
$$
\alpha:\Gamma_{\mathrm{BPS}}\longrightarrow\La^{2,1}_{II},
\qquad
\alpha(n,l,m)=2nf_2-lf_3+2mf_{-2}.
$$
Then
$$
(\alpha(\gamma),\alpha(\gamma'))
=-2\langle\gamma,\gamma'\rangle_{\mathrm{BPS}},
\qquad
(\alpha(\gamma),\alpha(\gamma))=-2(4nm-l^2).
$$
For $z=z_1f_{-2}+z_2f_3+z_3f_2$,
$$
q^nr^ls^m
=\exp(2\pi i(nz_1+lz_2+mz_3))
=\exp(-\pi i(\alpha(\gamma),z)).
$$
\end{proposition}

\begin{proof}
For $\gamma=(n,l,m)$ and $\gamma'=(n',l',m')$,
$$
(\alpha(\gamma),\alpha(\gamma'))
=2ll'-4nm'-4n'm
=-2\{2(nm'+n'm)-ll'\}.
$$
The character identity follows from
$$
(\alpha(n,l,m),z)=-2(nz_1+lz_2+mz_3).
$$
\end{proof}

\begin{remark}[What this proves physically]
The charge lattice, discriminant, and Lorentzian degree map are not
modular-form conventions.  They follow from the even two-charge Gram
matrix and from the chosen genus-two chemical potentials after passage
to the universal Gram-invariant index lattice.  What is not proved here
is the microscopic assertion that the full compactification index
factors through these three invariants for every physical sector of
$K3\times E$ curve counting, or that every formal triple in $\Z^3$ is
realized by an actual charge pair in the chosen compactification.
\end{remark}

\begin{proposition}[The duality group is mathematical; its physical
identification is input]
Let $\La^4=\Z e_1\oplus\cdots\oplus\Z e_4$, put
$$
q_0=e_1\wedge e_3+e_2\wedge e_4,
\qquad
\La^{3,2}=q_0^\perp\subset\wedge^2\La^4,
$$
and use the Pfaffian pairing on $\wedge^2\La^4$.  Then the
exterior-square representation gives
$$
\PSp_4(\Z)\simeq\SO_+(\La^{3,2})
\simeq\Orth(\La^{3,2})_+/\{\pm\Id_5\}.
$$
This identifies the Siegel period domain $\Hpl_2$ with the type-IV
domain of $\La^{3,2}$.
\end{proposition}

\begin{proof}
The subgroup of $\SL(\La^4)$ fixing $q_0$ under $\wedge^2$ is the
symplectic group of the alternating form
$$
B_{q_0}(x,y)\operatorname{vol}:=-x\wedge y\wedge q_0.
$$
The Pfaffian form on $\wedge^2\La^4$ has signature $(3,3)$, and
$q_0^\perp$ has signature $(3,2)$.  The exterior-square action preserves
$q_0^\perp$ and the positive type-IV component.  Its kernel on
$q_0^\perp$ is $\{\pm\Id_4\}$, giving the displayed exceptional
isomorphism.  The computation is the lattice calculation in the
exterior-square section; it is independent of the Borcherds product.
\end{proof}

\begin{remark}[Physical status of duality]
The preceding proposition is a theorem about integral lattices.  The
identification of this group with the physical duality group of the
compactification is an additional physical input.  The paper uses it
only in the following restricted sense: the protected determinant is an
automorphic section for this group, and the standard cusp is one
polarization of the same charge lattice.
\end{remark}

\begin{definition}[Active charge support]
Let
$$
\Gamma_{\mathrm{act}}
=\{\gamma=(n,l,m)\in\Gamma_{\mathrm{eff}}\mid f(nm,l)\neq0\}.
$$
By the weak-Jacobi support condition for $\phi_{0,1}\in J^{\mathrm{weak}}_{0,1}$,
$$
\gamma\in\Gamma_{\mathrm{act}}
\quad\Longrightarrow\quad
4nm-l^2\geq -1.
$$
The inactive triples in $\Gamma_{\mathrm{eff}}\setminus\Gamma_{\mathrm{act}}$
are ordering data only; their product exponents are zero.
\end{definition}

\begin{lemma}[Active support lands in the simple-root cone]
Let
$$
\delta_1=2f_2-f_3,\qquad
\delta_2=2f_{-2}-f_3,\qquad
\delta_3=f_3.
$$
For every $\gamma=(n,l,m)\in\Gamma_{\mathrm{act}}$,
$$
\alpha(\gamma)=n\delta_1+m\delta_2+(n+m-l)\delta_3,
$$
with
$$
n\geq0,\qquad m\geq0,\qquad n+m-l\geq0.
$$
Therefore
$$
\alpha(\Gamma_{\mathrm{act}})
\subset
\Z_{\geq0}\delta_1+\Z_{\geq0}\delta_2+\Z_{\geq0}\delta_3.
$$
\end{lemma}

\begin{proof}
The displayed decomposition is immediate:
$$
n(2f_2-f_3)+m(2f_{-2}-f_3)+(n+m-l)f_3
=2nf_2-lf_3+2mf_{-2}.
$$
The standard effective chamber gives $n,m\geq0$.  It remains to prove
$n+m-l\geq0$.  If $n=0$ or $m=0$, the support condition gives
$l^2\leq1$, and the three boundary cases in $\Gamma_{\mathrm{eff}}$
give $l\leq n+m$.  If $n,m>0$ and $l\geq n+m+1$, then
$$
l^2\geq(n+m+1)^2
=4nm+(n-m)^2+2(n+m)+1>4nm+1,
$$
contradicting $l^2\leq4nm+1$.  Hence $l\leq n+m$ in all cases.
\end{proof}

\begin{proposition}[Borcherds cone from active support]
The proof-grade Borcherds cone in the standard cusp is
$$
\Cone_{\mathrm{Borch}}
=\R_{\geq0}\,\alpha(\Gamma_{\mathrm{act}})
=\R_{\geq0}\delta_1+\R_{\geq0}\delta_2+\R_{\geq0}\delta_3.
$$
The three extremal rays are active:
$$
\delta_1=\alpha(1,1,0),\qquad
\delta_2=\alpha(0,1,1),\qquad
\delta_3=\alpha(0,-1,0),
$$
and each has discriminant $-1$ and square $2$.
\end{proposition}

\begin{proof}
The lemma gives the inclusion ``$\subset$''.  The three displayed
charges lie in $\Gamma_{\mathrm{eff}}$ and have
$$
f(0,1)=1,\qquad f(0,-1)=1,
$$
so they lie in $\Gamma_{\mathrm{act}}$.  Hence the three generators
$\delta_i$ lie in $\alpha(\Gamma_{\mathrm{act}})$, proving the reverse
inclusion of cones.  Since
$$
4nm-l^2=-1
$$
for the three displayed charges, the anti-isometry gives
$(\delta_i,\delta_i)=2$.
\end{proof}

\begin{remark}[Attack on the full-semigroup cone]
One must not define the Borcherds cone using all of
$\alpha(\Gamma_{\mathrm{eff}})$ unless zero-exponent triples are
removed.  For example
$$
\gamma=(0,2,1)\in\Gamma_{\mathrm{eff}},
\qquad
4nm-l^2=-4,
\qquad
f(0,2)=0,
$$
and
$$
\alpha(\gamma)=2f_{-2}-2f_3
=0\cdot\delta_1+1\cdot\delta_2-1\cdot\delta_3
\notin\R_{\geq0}\delta_1+\R_{\geq0}\delta_2+\R_{\geq0}\delta_3.
$$
Thus the chamber is forced by the active Jacobi support, not by the
lexicographic ordering alone.
\end{remark}

\begin{remark}[Residual physical assumption]
The derivation above proves the lattice algebra, the discriminant
pairing, the anti-isometry to $\La^{2,1}_{II}$, the exterior-square
duality group, and the active-support cone.  It does not prove from
first principles that the compactification produces the physical
duality group $\PSp_4(\Z)$, that every protected BPS index is a function
only of $(Q^2/2,Q\cdot P,P^2/2)$, that every formal Gram triple in
$\Gamma_{\mathrm{BPS}}$ is realized by a microscopic charge pair, or
that the automorphic real-root walls are microscopic walls of marginal
stability.  Those are the remaining physical assumptions.  This agrees
with the Vol.~III status: the Igusa boundary supplies the
decategorified denominator and positive half; the bracket-level
Hall--Drinfeld/BPS realization remains open.
\end{remark}
